\section{Procesor}\label{sec:hw-procesor}

Procesor vykonává strojové instrukce programu. Nejde však o~jedinou součástku provádějící všechny činnosti stejným způsobem. Uvnitř procesoru spolupracuje několik částí, z~nichž každá má jinou úlohu. Pro náš zjednodušený model budou nejdůležitější aritmeticko-logická jednotka, řadič a~registry.

\begin{itemize}
    \item \textbf{Aritmeticko-logická jednotka (ALU)} provádí aritmetické a~logické operace. Může například sčítat čísla, porovnávat hodnoty nebo provádět bitové operace.
    \item \textbf{Řadič} organizuje činnost procesoru. Na základě právě načtené instrukce vytváří řídicí signály, kterými určuje, odkud se mají získat operandy, jakou operaci má provést ALU a~kam se má uložit výsledek.
    \item \textbf{Registry} uchovávají malé množství právě používaných dat přímo uvnitř procesoru. Přístup k~nim je výrazně rychlejší než přístup do hlavní paměti.
\end{itemize}

Skutečný procesor obsahuje ještě řadu dalších částí, například jednotky pro práci s~desetinnými čísly, mechanismy pro předvídání větvení, několik úrovní cache nebo jednotky schopné provádět jednu operaci nad více hodnotami současně. Základní vztah mezi řadičem, výpočetními jednotkami, registry a~pamětí však zůstává užitečný i~při popisu podstatně složitějších procesorů.

\begin{figure}[ht]
    \centering
    \begin{tikzpicture}[
    x=1cm,
    y=1cm,
    >=Stealth,
    every node/.style={font=\small},
    unit/.style={
        draw,
        thick,
        rounded corners,
        align=center
    },
    subunit/.style={
        draw,
        rounded corners,
        fill=white,
        minimum width=2.2cm,
        minimum height=0.8cm,
        align=center
    },
    databus/.style={<->, very thick, blue!70!black},
    addrbus/.style={->, very thick, orange!85!black},
    ctrlbus/.style={<->, very thick, red!70!black}
]
    \node[
        unit,
        fill=blue!10,
        minimum width=6.0cm,
        minimum height=4.4cm
    ] (cpu) at (-3.1,1.1) {};
    \node[font=\large\bfseries] at (-4.8,2.7) {CPU};

    \node[subunit, fill=blue!18] (alu) at (-4.3,1.3) {\textbf{ALU}};
    \node[subunit, fill=orange!20] (control) at (-4.3,0.1) {\textbf{Řadič}};
    \node[
        subunit,
        fill=green!15,
        minimum width=2.5cm,
        minimum height=2.0cm
    ] (registers) at (-1.8,0.7) {\textbf{Registry}\\[1mm] $R_0\quad R_1$\\ $R_2\quad R_3$};

    \draw[->, thick] (registers.west) -- (alu.east);
    \draw[->, thick] (control.east) -- (registers.west);
    \draw[->, thick] (control.north) -- (alu.south);

    \node[
        unit,
        fill=green!12,
        minimum width=3.8cm,
        minimum height=2.7cm
    ] (memory) at (4.5,1.1) {\textbf{\large Paměť}\\[2mm] instrukce a~data};

    \draw[databus] (-5.6,-2.0) -- (6.4,-2.0);
    \draw[addrbus] (-5.6,-2.8) -- (6.4,-2.8);
    \draw[ctrlbus] (-5.6,-3.6) -- (6.4,-3.6);

    \node[anchor=east, text=blue!70!black] at (-5.8,-2.0) {datová};
    \node[anchor=east, text=orange!85!black] at (-5.8,-2.8) {adresová};
    \node[anchor=east, text=red!70!black] at (-5.8,-3.6) {řídicí};
    \node[anchor=west, font=\small\bfseries] at (6.55,-2.8) {sběrnice};

    \draw[databus] (-3.7,-1.1) -- (-3.7,-2.0);
    \draw[addrbus] (-3.1,-1.1) -- (-3.1,-2.8);
    \draw[ctrlbus] (-2.5,-1.1) -- (-2.5,-3.6);

    \draw[databus] (3.8,-0.25) -- (3.8,-2.0);
    \draw[addrbus] (4.5,-0.25) -- (4.5,-2.8);
    \draw[ctrlbus] (5.2,-0.25) -- (5.2,-3.6);
\end{tikzpicture}

    \caption{Zjednodušená vnitřní struktura procesoru a~jeho spojení s~pamětí.}
    \label{fig:hw-cpu-sbernice}
\end{figure}

\subsection{Sběrnice}\label{subsec:hw-sbernice}

Procesor, paměť a~další zařízení potřebují přenášet různé druhy informací. Soubor vodičů a~pravidel, podle nichž se po nich informace přenášejí, nazýváme \emph{sběrnice}. Na obrázku~\ref{fig:hw-cpu-sbernice} rozlišujeme tři logické skupiny signálů:
\begin{itemize}
    \item \textbf{datová sběrnice} přenáší samotná data a~instrukce;
    \item \textbf{adresová sběrnice} určuje paměťové místo nebo zařízení, s~nímž chce procesor komunikovat;
    \item \textbf{řídicí sběrnice} přenáší signály určující například směr přenosu, požadavek na čtení či zápis nebo informaci o~dokončení operace.
\end{itemize}

Uvažujme například čtení hodnoty z~paměti. Procesor nejprve předá na adresovou sběrnici adresu požadovaného místa a~řídicím signálem oznámí, že chce číst. Paměť vyhledá odpovídající obsah a~předá jej po datové sběrnici zpět procesoru. Při zápisu je postup podobný, pouze procesor zároveň odešle hodnotu a~řídicím signálem určí, že se má obsah paměti změnit.

\notebox{Rozdělení na datovou, adresovou a~řídicí sběrnici popisuje význam přenášených signálů. Ve skutečném zařízení nemusí každé skupině odpovídat samostatná sada vodičů; některé fyzické spoje mohou v~různých okamžicích přenášet různé druhy informací.}

\subsection{Instrukční cyklus}\label{subsec:hw-instrukcni-cyklus}

Program je pro procesor posloupností strojových instrukcí uložených v~paměti. Každá instrukce musí být načtena, rozpoznána a~provedena. Jednotlivé procesory se v~podrobnostech liší, ale základní průběh lze popsat následujícími kroky:
\begin{enumerate}
    \item \textbf{Čtení instrukce (\emph{fetch}).} Procesor získá z~paměti instrukci, na kterou ukazuje čítač instrukcí.
    \item \textbf{Dekódování instrukce (\emph{decode}).} Řadič určí, jakou operaci instrukce požaduje a~kde se nacházejí její operandy.
    \item \textbf{Čtení operandů (\emph{operand fetch}).} Potřebné hodnoty se získají z~registrů, případně z~paměti.
    \item \textbf{Provedení instrukce (\emph{execute}).} Příslušná výpočetní jednotka provede požadovanou operaci.
    \item \textbf{Uložení výsledku (\emph{write back}).} Výsledek se zapíše do registru nebo do paměti.
\end{enumerate}

\begin{figure}[ht]
    \centering
    \begin{tikzpicture}[
    x=1cm,
    y=1cm,
    >=Stealth,
    every node/.style={font=\small},
    stagebox/.style={
        draw,
        thick,
        rounded corners,
        fill=blue!12,
        align=center,
        text width=2.25cm,
        minimum height=1.25cm
    },
    arrow/.style={->, thick}
]
    \node[stagebox] (fetch) at (-5.6,0.8) {\textbf{Čtení}\\ instrukce\\[-1mm]\scriptsize fetch};
    \node[stagebox] (decode) at (-2.8,0.8) {\textbf{Dekódování}\\ instrukce\\[-1mm]\scriptsize decode};
    \node[stagebox] (operands) at (0,0.8) {\textbf{Čtení}\\ operandů\\[-1mm]\scriptsize operand fetch};
    \node[stagebox] (execute) at (2.8,0.8) {\textbf{Provedení}\\ instrukce\\[-1mm]\scriptsize execute};
    \node[stagebox] (writeback) at (5.6,0.8) {\textbf{Uložení}\\ výsledku\\[-1mm]\scriptsize write back};

    \draw[arrow] (fetch.east) -- (decode.west);
    \draw[arrow] (decode.east) -- (operands.west);
    \draw[arrow] (operands.east) -- (execute.west);
    \draw[arrow] (execute.east) -- (writeback.west);

    \node[
        draw,
        thick,
        rounded corners,
        fill=orange!18,
        align=center,
        minimum width=2.4cm,
        minimum height=0.9cm
    ] (interrupt) at (5.6,-1.6) {kontrola\\ přerušení};

    \draw[arrow] (writeback.south) -- (interrupt.north);
    \draw[arrow]
        (interrupt.west) -- (3.8,-1.6) -- (3.8,-2.5) --
        (-5.6,-2.5) -- (fetch.south);
\end{tikzpicture}

    \caption{Zjednodušený instrukční cyklus procesoru.}
    \label{fig:hw-instrukcni-cyklus}
\end{figure}

Představme si například instrukci, která má sečíst hodnoty registrů $R_0$ a~$R_1$ a~výsledek uložit do registru $R_2$. Procesor instrukci nejprve načte a~dekóduje. Poté přečte hodnoty obou zdrojových registrů, předá je ALU, nechá provést součet a~výsledek zapíše do cílového registru. Následně může pokračovat další instrukcí.

Obrázek~\ref{fig:hw-instrukcni-cyklus} vytváří dojem, že procesor dokončí všechny kroky jedné instrukce a~teprve potom začne načítat další. Takový postup je možný, moderní procesory však jednotlivé fáze často překrývají. Zatímco jedna instrukce se provádí, další se může dekódovat a~ještě další načítat. Tento způsob zpracování se nazývá \emph{zřetězení} neboli \emph{pipeline}. Překrývání fází zvyšuje počet instrukcí, které lze zpracovat za určitou dobu, přestože doba průchodu jedné instrukce všemi fázemi se nemusí zkrátit.

Činnost procesoru je řízena hodinovým signálem. Jeho frekvence udává počet taktů za sekundu, nikoli přímo počet dokončených instrukcí. Některá instrukce může vyžadovat více taktů, více instrukcí se může zpracovávat souběžně a~procesor může čekat na data z~paměti. Samotná hodnota frekvence proto nestačí k~porovnání výkonu různých procesorů.

\paragraph{Přerušení.}

Běžící program není jediným zdrojem událostí, na které musí procesor reagovat. Klávesnice může oznámit stisk klávesy, úložiště dokončení přenosu nebo časovač uplynutí stanoveného intervalu. K~takovým situacím slouží \emph{přerušení}. Procesor po dokončení vhodného kroku uloží informace potřebné pro pozdější pokračování programu a~přejde k~obsluze vzniklé události. Po jejím vyřízení se může k~původnímu programu vrátit.

Přerušení umožňuje, aby procesor nemusel každé zařízení neustále kontrolovat. Zařízení jej samo upozorní ve chvíli, kdy potřebuje obsluhu. Operační systém pak rozhodne, jaká část programu má událost zpracovat a~kdy může pokračovat přerušený výpočet.

\subsection{Registry}\label{subsec:hw-registry}

Registry jsou velmi malé a~rychlé paměti přímo uvnitř procesoru. Uchovávají operandy právě prováděných instrukcí, adresy, mezivýsledky i~informace o~stavu výpočtu. Instrukce obvykle pracují s~registry přímo, zatímco hodnotu z~hlavní paměti je nejprve potřeba do registru načíst a~výsledek případně později zapsat zpět.

\begin{figure}[ht]
    \centering
    \begin{tikzpicture}[
    x=1cm,
    y=1cm,
    >=Stealth,
    every node/.style={font=\small},
    block/.style={
        draw,
        thick,
        rounded corners,
        align=center
    },
    reg/.style={
        draw,
        thick,
        rounded corners,
        fill=green!14,
        minimum width=1.0cm,
        minimum height=0.65cm
    },
    special/.style={
        draw,
        thick,
        rounded corners,
        fill=orange!20,
        minimum width=1.65cm,
        minimum height=0.65cm,
        font=\small\bfseries
    }
]
    \node[
        block,
        fill=blue!8,
        minimum width=13.5cm,
        minimum height=5.1cm
    ] (cpu) at (0,0) {};
    \node[font=\large\bfseries] at (-5.2,1.85) {Procesor};

    \node[
        block,
        fill=blue!18,
        minimum width=2.4cm,
        minimum height=0.95cm
    ] (alu) at (-4.7,0.55) {\textbf{ALU}};
    \node[
        block,
        fill=orange!22,
        minimum width=2.4cm,
        minimum height=0.95cm
    ] (control) at (-4.7,-0.75) {\textbf{řadič}};

    \node[
        block,
        fill=green!7,
        minimum width=5.1cm,
        minimum height=3.3cm
    ] (regfile) at (0,0) {};
    \node[font=\bfseries] at (0,1.15) {pracovní registry};
    \node[reg] at (-1.45,0.35) {$R_0$};
    \node[reg] at (0,0.35) {$R_1$};
    \node[reg] at (1.45,0.35) {$R_2$};
    \node[reg] at (-1.45,-0.65) {$R_3$};
    \node[reg] at (0,-0.65) {$R_4$};
    \node[reg] at (1.45,-0.65) {$R_5$};

    \node[special] (pc) at (4.8,1.35) {PC};
    \node[special] (ir) at (4.8,0.45) {IR};
    \node[special] (sp) at (4.8,-0.45) {SP};
    \node[special] (flags) at (4.8,-1.35) {FLAGS};

    \draw[->, thick] (alu.east) -- (regfile.west |- alu.east);
    \draw[->, thick] (control.east) -- (regfile.west |- control.east);
    \draw[->, thick] (ir.west) -- (regfile.east |- ir.west);
\end{tikzpicture}

    \caption{Pracovní a~vybrané speciální registry v~procesoru.}
    \label{fig:hw-registry}
\end{figure}

Vedle pracovních registrů, jejichž význam určuje právě prováděný program, se setkáme také s~registry se zvláštní úlohou:
\begin{itemize}
    \item \textbf{PC (\emph{program counter})} obsahuje adresu následující instrukce, která se má načíst. Po běžné instrukci se posune vpřed, zatímco instrukce skoku do něj může uložit jinou adresu.
    \item \textbf{IR (\emph{instruction register})} uchovává právě načtenou nebo prováděnou instrukci.
    \item \textbf{SP (\emph{stack pointer})} ukazuje na vrchol zásobníku v~paměti. Zásobník se využívá například při volání funkcí.
    \item \textbf{FLAGS} uchovává příznaky výsledku poslední operace, například zda byl výsledek nulový, záporný nebo zda došlo k~přetečení.
\end{itemize}

Konkrétní názvy, počty a~přesný význam registrů závisí na instrukční sadě procesoru. Obrázek~\ref{fig:hw-registry} proto nepředstavuje schéma jednoho konkrétního modelu, ale obecnou ilustraci rolí, které registry mohou zastávat.

\importantbox{Registry nejsou určeny k~dlouhodobému uložení velkého množství dat. Jejich hlavní výhodou je rychlost a~přímá dostupnost pro instrukce procesoru; za tuto výhodu platíme velmi malou kapacitou.}
